\chapter{Learning with an explicit error budget before a dependent evaluation}
\label{exec13:dependent-pilot}

Consider an encouragement to take a bank offer. The bank observes information
$X$ before the encouragement $Z$, the eventual action $A$, and its subsequent
value $Y$. Flexible prediction is useful because accounts differ in their
pre-existing circumstances. It also creates a second estimation problem:
removing predictable variation imperfectly changes the instrument moment.
When the residual response to the encouragement is small, that change can be
large relative to the information identifying the effect. A narrow interval
computed after fitting a good predictor need not describe the uncertainty in
the bank's causal decision.

This chapter develops a restricted alternative. Learn on an independent pilot,
obtain an explicit bound on each prediction error, and carry that bound into
a joint region for two observable moments. Evaluate those moments on a bounded
sequence with a known dependence range. Project the region into the effect,
retaining disconnected and unbounded answers. The construction is a project
derivation with exact numerical references, rather than an implemented general
learner or a new repeated-sampling qualification~\citep{exec13_dependent_argument}.
Its restrictions are substantial; explaining them is part of deciding whether
the proposed product merits a larger experiment.

\section{What the repaired predecessor did and did not establish}

The preceding sparse-learning experiment already has a repaired version. Its
first version selected penalties using evaluation outcomes, standardized with
all observations, included covariance products across an omitted calendar gap,
and ignored some scenario dimensions. The repaired version uses training-only
nested selection, training-specific scaling, within-fold Bartlett covariance,
and the actual registered dimension and support. The original results remain
exploratory. They do not establish the performance of that repaired procedure.

For the repaired learner, the high-dimension and larger-support scenarios
covered respectively $499/512$ and $493/512$ requests. Their simultaneous lower
coverage bounds, $0.949140$ and $0.933851$, fell below the frozen $0.95$ threshold.
The nonlinear omitted-nuisance scenario covered $0/512$. The weak-instrument
scenario covered $512/512$, but its median interval width was $3.836470$ in the
synthetic effect units. These were nominal $99\%$ intervals, so comparisons
with the original $95\%$ intervals cannot isolate the effect of repairing the
code. Failure to pass a certification screen is distinct from proof of universal
undercoverage; complete finite availability is also distinct from useful
precision. These results remain the relevant negative baseline.

An independent pilot removes one source of dependence between fitted functions
and evaluation scores. It cannot make a misspecified dictionary correct, turn
a weak instrument into a strong one, or make repeated calendar observations
independent. A successor therefore needs a statement for each of those issues.

\section{The structural target and its observable moment}

The imposed partially linear model is
\begin{equation}
 Y=\theta_0 A+g_0(X)+\epsilon,
 \qquad E[\epsilon\mid X,Z]=0.
 \label{exec13:pilot-pliv}
\end{equation}
The conditional moment excludes a remaining direct effect of the encouragement
and unmeasured causes associated with its residual after conditioning on $X$.
It is a scientific assumption. A prediction fit or source hash cannot verify it.
Here $\theta_0$ is a constant structural coefficient, expressed in outcome
units per unit of action. Interpreting it as an unrestricted heterogeneous
average effect would require a different argument.

Let $m(X)=E[Z\mid X]$, $r(X)=E[A\mid X]$ and
$\ell(X)=E[Y\mid X]$. Define residuals
$v_Z=Z-m(X)$, $v_A=A-r(X)$ and $v_Y=Y-\ell(X)$, and moments
\begin{equation}
 M=E[v_Zv_Y],\qquad Q=E[v_Zv_A].
 \label{exec13:pilot-moments}
\end{equation}
Equation~\eqref{exec13:pilot-pliv} implies $M=\theta_0 Q$. A nonzero $Q$
identifies the scalar coefficient from this relation. If $Q=0$, the relation
can be uninformative even when the observations themselves are measured exactly.

The residualized score
$(Z-m(X))\{Y-\ell(X)-\theta(A-r(X))\}$ is the inherited partially linear IV
score in \citet[Section~4.2, equation~(4.8)]{dml2018}. That source states its
moment, nuisance-product and nonzero-slope conditions. Its independent-sample
theorem does not supply a result for arbitrary serially dependent bank rows.
The following finite construction instead proves its own statement for a
declared observation law.

\section{An observation law that makes separation meaningful}

The pilot contains $m_P$ independent draws from the marginal law of
$W=(X,Z,A,Y)$ and is independent of the entire evaluation sequence. The
evaluation observations $W_1,\ldots,W_n$ are stationary and $d$-dependent:
collections separated by more than $d$ dates generate independent sigma fields.
The integer $d$ is supplied and justified by the design. Uncorrelated residuals
or a small sample autocorrelation are insufficient to establish this property.

Assume known almost-sure ranges $|Z|\le B_Z$, $|A|\le B_A$, and $|Y|\le B_Y$.
Clip predicted conditional means into the respective ranges. Since the true
conditional means lie there, clipping predictions cannot increase their squared
error. Clipping the observed outcome is different: it can change the target
and cannot be introduced silently to make a bound available.

The features $\phi(X)\in\mathbb R^p$ are fixed before the pilot, with
$|\phi_k(X)|\le K$. A feature learned from data needs a separate design sample
and conditional versions of these assumptions. A constant feature can represent
an intercept, but the derivation below penalizes its coefficient too. It does
not silently inherit the previous implementation's unpenalized intercept or
its data-dependent scaling.

These choices exclude the unbounded Gaussian AR experiment in the predecessor.
They can describe a bounded generated benchmark or a justified collection of
finite-dependent observations. They do not describe an arbitrary bank panel.
An extension to geometric mixing requires a proved coupling or concentration
bound, with an explicit additional error allocation. Unknown dependence remains
an information requirement.

\section{The nuisance error enters as two products}

Condition on the independent pilot, so the fitted functions are fixed. Write
$\delta_Z=\widehat m-m$, $\delta_A=\widehat r-r$ and
$\delta_Y=\widehat\ell-\ell$. Expanding the product, then conditioning on $X$,
eliminates the linear error terms because $E[v_j\mid X]=0$. Thus
\begin{align}
 M_{\rm fit}&=E[(Z-\widehat m)(Y-\widehat\ell)\mid\text{pilot}]
             =M+E[\delta_Z\delta_Y\mid\text{pilot}],\\
 Q_{\rm fit}&=E[(Z-\widehat m)(A-\widehat r)\mid\text{pilot}]
             =Q+E[\delta_Z\delta_A\mid\text{pilot}].
 \label{exec13:pilot-bias}
\end{align}
If a simultaneous certificate supplies $\|\delta_j\|_2\le e_j$, Cauchy--Schwarz
gives
\begin{equation}
 |M_{\rm fit}-M|\le e_Ze_Y,\qquad
 |Q_{\rm fit}-Q|\le e_Ze_A.
 \label{exec13:pilot-products}
\end{equation}
Consequently the score bias at a candidate $\theta$ is bounded by
$e_Z(e_Y+|\theta|e_A)$. The encouragement prediction error is essential to
this expression. An outcome--action product alone does not describe the IV
score. Nor does an error norm of exact order $n^{-1/4}$ automatically suffice:
two such errors can leave a root-$n$ scaled remainder of order one.

The bound is exact in its conditional decomposition and conservative in its
norm inequality. It needs neither an estimated effect in the denominator nor
a claim that the remainder is negligible. Those features allow a finite
construction that remains honest when a nuisance certificate is wide.

\section{A sparse learner with a stated population-error certificate}

For each response $T_j\in\{Z,A,Y\}$, suppose its conditional mean has the
representation
\begin{equation}
 f_j(X)=\phi(X)^\top\beta_j+a_j(X),\qquad
 |\operatorname{supp}(\beta_j)|\le s_j,\qquad |a_j(X)|\le A_j.
\end{equation}
Let $\epsilon_j=T_j-f_j(X)$ satisfy
$E[\epsilon_j\mid X]=0$ and $|\epsilon_j|\le E_j$. The approximation bound
$A_j$ is a premise; it is not a training residual. In the exact sparse model
$A_j=0$. Let $\Phi$ be the pilot feature matrix, and fit the exact minimizer
\begin{equation}
 \widehat\beta_j\in\arg\min_b
 \left\{\frac{\|T_j-\Phi b\|_2^2}{2m_P}+\lambda_j\|b\|_1\right\}.
 \label{exec13:pilot-lasso}
\end{equation}

To see where the penalty requirement comes from, each coordinate of
$\phi(X)\epsilon_j$ is centered and bounded in an interval of width $2KE_j$.
Its exponential moment is at most
$\exp\{t^2(2KE_j)^2/8\}$. Independence, the exponential Markov inequality,
optimization in $t$ and a union over $p$ coordinates give, with error at most
$\eta_j$,
\begin{equation}
 \|\Phi^\top\epsilon_j/m_P\|_\infty
 \le KE_j\sqrt{\frac{2\log(2p/\eta_j)}{m_P}}.
\end{equation}
The approximation residual contributes at most $KA_j$. A sufficient penalty is
\begin{equation}
 \lambda_j\ge2K\left\{E_j\sqrt{\frac{2\log(2p/\eta_j)}{m_P}}+A_j\right\}.
 \label{exec13:pilot-penalty}
\end{equation}
A fixed finite family whose every member meets this lower bound shares the
same event. A rate statement also requires that the selected penalty is at most
a fixed multiple of the threshold. An arbitrary validation-selected penalty
does not receive that property merely because its predictive loss is small.

For $h=\widehat\beta_j-\beta_j$ and $S=\operatorname{supp}(\beta_j)$,
optimality in \eqref{exec13:pilot-lasso} and the preceding score event yield
\begin{equation}
 \frac{\|\Phi h\|_2^2}{m_P}+\lambda_j\|h_{S^c}\|_1
 \le3\lambda_j\|h_S\|_1.
 \label{exec13:pilot-basic}
\end{equation}
In particular, $h$ belongs to the cone
$\|h_{S^c}\|_1\le3\|h_S\|_1$. Suppose a certificate establishes
\begin{equation}
 \frac{\|\Phi h\|_2^2}{m_P}\ge\kappa_j^2\|h_S\|_2^2
 \label{exec13:pilot-cone}
\end{equation}
for that cone and every admissible support of size at most $s_j$. Combining
the two inequalities gives
\begin{equation}
 \frac{\|\Phi h\|_2}{\sqrt{m_P}}\le
 \frac{3\sqrt{s_j}\lambda_j}{\kappa_j},\qquad
 \|h\|_1\le\frac{12s_j\lambda_j}{\kappa_j^2}.
\end{equation}
This is an empirical prediction bound. It still needs transfer to the population
in which the evaluation takes place.

Write $\widehat\Sigma=\Phi^\top\Phi/m_P$ and
$\Sigma=E[\phi(X)\phi(X)^\top]$. The same bounded-mean argument, applied to
all $p^2$ feature products, gives with error $\eta_G$
\begin{equation}
 \|\widehat\Sigma-\Sigma\|_{\max}\le
 \tau=K^2\sqrt{\frac{2\log(2p^2/\eta_G)}{m_P}}.
\end{equation}
Since $|h^\top(\Sigma-\widehat\Sigma)h|\le\tau\|h\|_1^2$, a usable
population certificate is
\begin{equation}
 e_j=A_j+
 \sqrt{\frac{9s_j\lambda_j^2}{\kappa_j^2}
       +\frac{144\tau s_j^2\lambda_j^2}{\kappa_j^4}}.
 \label{exec13:pilot-certificate}
\end{equation}
Prediction clipping preserves this bound. The three score events and the Gram
event have total failure probability at most
$\eta_P=\eta_Z+\eta_A+\eta_Y+\eta_G$. They may use the same pilot;
the union bound does not require them to be independent.

Under exact approximation, bounded constants, a positive cone eigenvalue, and
bounded $\tau s_j/\kappa_j^2$, the prediction norm is of order
$\sqrt{s_j\log(p)/m_P}$. The squared error, rather than the norm, has order
$s_j\log(p)/m_P$. For comparable nuisance complexities, negligible root-$n$
bias additionally requires
$\sqrt n\,s\log(p)/m_P\longrightarrow0$. The finite construction below
retains the explicit bias bound even when this asymptotic requirement is not
plausible at the available sample sizes.

\subsection{What must actually be certified}

The cone bound is a substantive requirement. A full Gram condition number is
not a sparse cone certificate when $p$ exceeds the pilot size. A small fixed
dictionary can use a certified minimum eigenvalue. A high-dimensional
construction needs a justified sparse certificate or a generated-design
guarantee. If certification fails, the result must declare the bound unavailable
or use an explicitly wider fallback. The theorem cannot assume that a returned
regression coefficient certifies its own rate.

The minimization also needs a numerical specification. An approximate optimizer
adds a basic-inequality error that must be propagated. A small KKT residual
does not by itself prove equation~\eqref{exec13:pilot-certificate}. The exact
reference uses a 32-row, four-feature orthogonal design, support two and penalty
$0.1$. Soft thresholding gives coefficients $(0.4,0.15,0,0)$; a separate
nonnegative-variable convex optimization agrees. Population squared error is
$0.02$, below the selected bound $0.18$. This checks the algebra and the small
optimization reference; it supplies no certificate for an arbitrary pilot.

\subsection{Why a proved bound may still be economically unattractive}

Table~\ref{exec13:pilot-cost-table} evaluates the formula under the hypothetical
certified constants $K=E_j=\kappa_j=1$, $A_j=0$, and error allocation $0.005$
for each score and Gram event. These are planning inputs, not fitted properties
of the predecessor. Even favorable assumed geometry gives a loose bound at
small pilot sizes. A deterministic range cap could reduce the largest numbers,
but would not make them a useful precision claim.

\begin{table}[htbp]
\centering\small
\caption{Uncapped nuisance norm and product certificates under stated hypothetical constants.}
\label{exec13:pilot-cost-table}
\begin{tabular}{rrrrr}
\toprule
Features & Support & Pilot rows & Norm bound & Equal-error product\\
\midrule
20 & 3 & 1,024 & 3.9770 & 15.8166\\
60 & 3 & 1,024 & 4.3729 & 19.1221\\
20 & 8 & 1,024 & 10.2004 & 104.0480\\
60 & 8 & 65,536 & 0.5698 & 0.3247\\
60 & 8 & 1,048,576 & 0.0961 & 0.0092\\
\bottomrule
\end{tabular}
\end{table}

The table does not recommend collecting a million bank observations. It shows
why a prototype should compare sharper justified bounds, smaller dictionaries,
better information and direct calibration before proposing an expensive bank
experiment. A valid but uninformative theorem is progress in defining the
problem, not evidence that the product already delivers an economical answer.

\section{Independent calibration as a different information route}

When sparse structure cannot be certified, a frozen clipped learner can be
evaluated on an additional independent calibration sample of $k$ marginal draws.
For response $T_j$, let $L_j$ be average squared predictive loss and let
$v_j^-$ be a justified lower bound on $E[\epsilon_j^2]$. Conditional mean
orthogonality gives
\begin{equation}
 E[(T_j-\widehat f_j(X))^2\mid\text{pilot}]
 =\|\widehat f_j-f_j\|_2^2+E[\epsilon_j^2].
\end{equation}
The observed loss is in $[0,4B_j^2]$. Therefore, with error $\eta_{Cj}$,
\begin{equation}
 \|\widehat f_j-f_j\|_2^2\le
 L_j+4B_j^2\sqrt{\frac{\log(1/\eta_{Cj})}{2k}}-v_j^-.
 \label{exec13:pilot-calibration}
\end{equation}
Use the nonnegative part before taking a square root and cap the norm by the
deterministic $2B_j$ bound. A negative uncapped expression is a diagnostic for
an inconsistent noise bound or a rare confidence event; it is not proof of
perfect learning. Joint error accounting is required if the analyst plans to
select the tighter of multiple certificates after observing them.

The safe choice $v_j^-=0$ needs no noise decomposition. It can leave an
irreducible-noise floor that prevents the certificate from shrinking. Better
noise information can come from a known generated mechanism, independently
justified repeated measurements, or an additional model with its own audit.
This route exchanges structural learning assumptions for calibration data and
noise information. Its full cost includes both samples.

\section{A finite region under a known dependence range}

Set $q=d+1\le n$ and split the evaluation dates by residue modulo $q$. Within
each class the observations are independent, conditional on the independent
pilot. Different classes can remain dependent. The fitted products
$U_{Mt}=(Z_t-\widehat m(X_t))(Y_t-\widehat\ell(X_t))$ and
$U_{Qt}=(Z_t-\widehat m(X_t))(A_t-\widehat r(X_t))$ lie in intervals of
widths $8B_ZB_Y$ and $8B_ZB_A$ respectively.

Let $n_{\min}=\lfloor n/q\rfloor$. A two-sided bounded-mean inequality for
each of the $2q$ class/coordinate combinations and a union bound give
\begin{equation}
 t_j=8B_ZB_j
 \sqrt{\frac{\log(4q/\alpha_E)}{2n_{\min}}},
 \qquad j\in\{Y,A\}.
 \label{exec13:pilot-temporal-band}
\end{equation}
With probability at least $1-\alpha_E$, both overall sample means deviate from
their fitted population moments by at most their respective $t_j$. To verify
the final step, the overall mean is a weighted average of the class means;
bounding every class bounds that weighted average even when class sizes differ.
Neither independence between classes nor independence between numerator and
denominator is assumed.

Let $\widehat M=n^{-1}\sum_tU_{Mt}$ and
$\widehat Q=n^{-1}\sum_tU_{Qt}$. With
$w_M=t_Y+e_Ze_Y$ and $w_Q=t_A+e_Ze_A$, form
\begin{equation}
 C=[\widehat M-w_M,\widehat M+w_M]
   \mathbin\times[\widehat Q-w_Q,\widehat Q+w_Q].
 \label{exec13:pilot-rectangle}
\end{equation}
For the sparse route, $C$ covers $(M,Q)$ with probability at least
$1-\eta_P-\alpha_E$, provided the supplied certificates and model premises
hold. The calibration route replaces $\eta_P$ by the sum of its calibration
errors. The coverage statement concerns the joint rectangle, so later projection
needs no further multiplicity adjustment.

If the nuisance certificate is missing, $e_j=2B_j$ is an unconditional bound
for clipped predictions and bounded conditional means, but it may make $C$
useless. Alternatively the product can withhold a qualified set and report
which information is missing. It must not treat an empirical fit score as the
missing certificate. Likewise, the assertion $W_t=W_1$ for every date describes
a bounded process with arbitrarily many rows but constant mean variance.
Row count alone gives no universal $n^{-1/2}$ precision improvement.

\paragraph{Availability and selection.}
If a pilot fails its cone or numerical certificate, returning the whole line
preserves the all-request coverage statement. Withholding the result instead
is an availability failure and must stay in the denominator. The argument
bounds false exclusions across all requests; it does not establish the same
conditional coverage among only those requests whose certificates pass, or a
minimum fraction of informative results.

\section{Weak-instrument sets and a decision threshold}

Define
\begin{equation}
 S(C)=\{\theta:\exists(M,Q)\in C\text{ with }M=\theta Q\}.
 \label{exec13:pilot-projection}
\end{equation}
On the event that the true moments lie in $C$, their structural relation puts
$\theta_0$ in $S(C)$. This proves coverage without dividing by $\widehat Q$
or treating an estimated denominator as known. It is a set based on this
moment relation and the imposed model; it is not claimed sharp for the entire
observed distribution.

For a proposed value $\theta$, multiply the denominator interval by $\theta$
and test its intersection with the numerator interval, reversing endpoints when
$\theta<0$. This gives a direct feasibility reference independent of a ratio
interval implementation. Exact rational arithmetic checked 18,685 membership
comparisons across 225 rectangles, including boundary points. Examples are
\begin{align*}
 [1,2]\mathbin\times[-1/2,1/2]&\quad\longmapsto\quad
 (-\infty,-2]\cup[2,\infty),\\
 [-1,1]\mathbin\times[-1/2,1/2]&\quad\longmapsto\quad\mathbb R,\\
 [1,2]\mathbin\times\{0\}&\quad\longmapsto\quad\varnothing,\\
 \{1\}\mathbin\times\{2\}&\quad\longmapsto\quad\{1/2\}.
\end{align*}
An empty set can reject the imposed moment relation or reflect a confidence
failure. It does not mean a zero effect. An unbounded set can correctly express
inadequate information and still fail to support a banking action.

For a positive denominator interval and nonnegative numerator interval, the
projection is
\begin{equation}
 \left[\frac{\widehat M-w_M}{\widehat Q+w_Q},
       \frac{\widehat M+w_M}{\widehat Q-w_Q}\right],
 \qquad \widehat Q>w_Q,
\end{equation}
with width
\begin{equation}
 \frac{2(\widehat Qw_M+\widehat Mw_Q)}{\widehat Q^2-w_Q^2}.
 \label{exec13:pilot-width}
\end{equation}
Width diverges as the lower denominator approaches zero. The illustrative
rectangle $[0.91,1.01]\times[0.78,0.82]$ gives
$[91/82,101/78]\simeq[1.109756,1.294872]$. Its radii are hypothetical
inputs, not fitted pilot evidence. A sign-normalized action cost of $1.15$
falls inside this interval, leaving the action unresolved. More independent
training can shrink the nuisance products; more effective evaluation data can
shrink the temporal bands. Neither repairs a zero instrument or establishes
the validity of an exclusion assumption.

\section{An exact explanation for the nonlinear failure}

The repaired predecessor adds $h(X)=0.6(X_1^2-1)$ to the encouragement and
outcome equations. For a standard Gaussian coordinate, $E[h]=0$ and
$E[h^2]=0.72$; the term is orthogonal to the linear Gaussian features.
Even with unlimited pilot data, a linear dictionary therefore leaves $h$
in the encouragement residual and leaves $\theta_0kh+h$ in the outcome
residual, where $k$ is the encouragement coefficient in the action equation.
With unit encouragement-innovation variance, the residualized numerator and
denominator obey
\begin{equation}
 M_{\rm linear}=\theta_0 k(1+0.72)+0.72,\qquad
 Q_{\rm linear}=k(1+0.72).
\end{equation}
For $k=0.8$, the ratio converges to
\begin{equation}
 \theta_0+\frac{0.72}{0.8(1+0.72)}
 =\theta_0+\frac{45}{86}
 \simeq\theta_0+0.523255814.
\end{equation}
This is a population approximation-error explanation, not an exact prediction
of the saved finite-sample lasso bias. It explains why more observations alone
need not cure the failure. Prespecifying the quadratic feature removes this
particular population omission, but finite estimation and dependence still
need their own analysis. Under this Gaussian example the omitted feature is
unbounded, so the earlier bounded approximation theorem does not apply unchanged.

\section{The proposed next experiment and its limits}

A useful bounded prototype would compare a small fixed dictionary, the
certified sparse pilot, and the independent calibration route on the same
declared target. Its full resource record must include feature preparation,
training, tuning, certificate construction, calibration, dependent evaluation,
set projection and retained refusals. A faster optimizer is of little value if
the required pilot or confidence width makes the decision uneconomical.

Before final draws, choose a bounded mechanism with an explicit dependence
range and a separate pilot, freeze the dictionary and approximation bounds,
review the cone and numerical certificates, and profile complete costs. Then
freeze sample sizes, error allocations, adverse cases, alternatives, fresh
draws, stopping rule and useful-width criterion. The old failed outcomes cannot
become the new evaluation data. No final successor study is registered by the
argument in this chapter.

The route is now explicit for this restricted family. General unbounded
dependent learning, unknown approximation error, unverified instruments,
arbitrary learners and useful weak-IV precision remain open. Where the
certificate is too costly or too wide, narrowing the question or changing the
information design is a substantive product recommendation.
\label{exec13:dependent-pilot-end}
